% corpus_supplement: Erdős problem #287
% source: https://www.erdosproblems.com/latex/287
% fetched_at: 2026-07-14T18:48:13Z
% payload_sha256: 269fa40cd1c77aec2a0533d39b93ce88c0efbb3ac7ee59d04c5e448e88cc3ce7
% statement/remarks/references extracted by erdos_ingest.extract_source_page;
% rendered by erdos_ingest.render_tex — do not edit
\documentclass{article}
\usepackage{amsmath, amssymb, amsthm}
\usepackage{hyperref}
\usepackage{geometry}
\geometry{margin=1in}

\title{Erd\H{o}s Problem \#287}
\date{Last updated: 2025-12-5}
\author{Source: erdosproblems.com}

\begin{document}
\maketitle

\noindent\textbf{Status:} OPEN \quad \textbf{No prize}

\noindent\textbf{Tags:} number theory, unit fractions

\medskip
\noindent\textbf{Problem Statement:}

Let $k\geq 2$. Is it true that, for any distinct integers $1<n_1<\cdots <n_k$ such that\[1=\frac{1}{n_1}+\cdots+\frac{1}{n_k}\]we must have $\max(n_{i+1}-n_i)\geq 3$?

\bigskip
\noindent\textbf{Remarks:}

The example $1=\frac{1}{2}+\frac{1}{3}+\frac{1}{6}$ shows that $3$ would be best possible here. The lower bound of $\geq 2$ is equivalent to saying that $1$ is not the sum of reciprocals of consecutive integers, proved by Erd\H{o}s \cite{Er32}. 

This conjecture would follow for all but at most finitely many exceptions if it were known that, for all large $N$, there exists a prime $p\in [N,2N]$ such that $\frac{p+1}{2}$ is also prime.

\bigskip
\noindent\textbf{References:}

[Er32] Erd\H{o}s, P., Egy K\"{u}rschak-f\'{e}le elemi sz\'{a}melm\'{e}ti t\'{e}tel \'{a}ltad\'{a}nosit\'{a}sa. MAt. es Phys. Lapok (1932).

\bigskip
\noindent\small{Source: \url{https://www.erdosproblems.com/287}}
\end{document}
